\documentclass[dvipdfmx]{jsarticle}

\usepackage{ascmac}
\usepackage{url}
\usepackage[dvipdfmx]{hyperref}
\usepackage{pxjahyper}
\usepackage[dvipdfmx]{graphicx}
\usepackage{float}
\usepackage{listings,jlisting}

\hypersetup{
  colorlinks=true,
  urlcolor=cyan,
  linkcolor=black
}

\lstset{
  basicstyle={\ttfamily},
  identifierstyle={\small},
  commentstyle={\smallitshape},
  keywordstyle={\small\bfseries},
  ndkeywordstyle={\small},
  stringstyle={\small\ttfamily},
  frame={tb},
  breaklines=true,
  columns=[l]{fullflexible},
  numbers=left,
  xrightmargin=0zw,
  xleftmargin=3zw,
  numberstyle={\scriptsize},
  stepnumber=1,
  numbersep=1zw,
  lineskip=-0.5ex
}


\begin{document}

\section{実験目的・課題}

次の画像(図\ref{fig:tapu})について、以下の3つの処理を行う。

\begin{itemize}
  \item 画像に対して輪郭検出を行う
  \item 画像に対してフィルタ処理を行う
  \item 画像に対して自動検出によるマスク処理を行う
\end{itemize}

\begin{figure}[H]
  \centering
  \includegraphics[width=0.7\hsize]{../pic/tapu.png}
  \caption{画像処理を行う画像}
  \label{fig:tapu}
\end{figure}


\section{実装方法}

\subsection{画像の二値化}

画像の二値化とは、画素の特定の要素において閾値を与え、
0か1に分類することである。
二値化された画像は、図\ref{fig:binarization}のように白黒画像なる。
\begin{figure}[htbp]
  \centering
  \includegraphics[width=0.7\hsize]{../pic/binarization.jpg}
  \caption{大津の二値化を適用した図\ref{fig:tapu}}
  \label{fig:binarization}
\end{figure}

pythonのOpenCVではthresholdという関数で画像の二値化を行うことが出来る。
第1引数に入力画像、第2引数に閾値、第3引数は閾値以上の値を持つ画素に
割り当てられる値、第4引数にいくつかある閾値処理のどれを使用するかのフラグを受け取る。
返り値は2つあり、使用した閾値と、二値化を適用した画像である。
入力画像はグレースケール画像でなければならず、画像を読み込むときに
imread(name,0)のように第2引数に0を与えると図\ref{fig:grayscale}のように
グレースケール画像として読み込むことが出来る。
\begin{figure}[htbp]
  \centering
  \includegraphics[width=0.7\hsize]{../pic/grayscale.png}
  \caption{グレースケール画像}
  \label{fig:grayscale}
\end{figure}

\newpage
\subsection{輪郭検出}

輪郭とは、同じ色などについて境界に沿った連続する点をつなげることによって
形成される曲線のことである。
精度よく輪郭を検出するために、二値化画像を使用する。
以下のコードで輪郭を検出し、元画像に描画することが出来る。
\begin{lstlisting}
  import cv2
  img = cv2.imread('tapu.jpg')
  gray = cv2.cvtColor(img, cv2.COLOR_BGR2GRAY)
  ret, img_thres = cv2.threshold(gray, 0, 255, cv2.THRESH_OTSU)
  img_cont, contours, hierarchy = cv2.findContours(img_thres, cv2.RETR_EXTERNAL, cv2.CHAIN_APPROX_SIMPLE)
  cv2.drawContours(img, contours, -1, color=(0, 0, 255), thickness=2)
\end{lstlisting}

findContoursの第二引数には、全ての輪郭を同等に扱うRETR\_LIST、
入れ子構造になっている輪郭のもっとも外側の輪郭のみを検出するRETR\_EXTERNALなどを指定する。
第三引数では、全ての点を検出するCHAIN\_APPROX\_NONE、冗長な点を削除するCHAIN\_APPROX\_SIMPLE等を指定する。

\begin{figure}[htbp]
  \begin{minipage}{0.5\hsize}
    \begin{center}
      \includegraphics[width=0.9\hsize]{../pic/contour1_list.png}
    \end{center}
    \caption{RETR\_LISTで指定した場合}
    \label{fig:contour_list}
  \end{minipage}
  \begin{minipage}{0.5\hsize}
    \begin{center}
      \includegraphics[width=0.9\hsize]{../pic/contour1_external.png}
    \end{center}
    \caption{RETR\_EXTERNALで指定した場合}
    \label{fig:contour_external}
  \end{minipage}
\end{figure}

図\ref{fig:contour_list}に、findContoursの第二引数にRETR\_LISTを指定したもの、
図\ref{fig:contour_external}に、RETR\_EXTERNALを指定したものを示す。
図\ref{fig:contour_external}では、図\ref{fig:contour_list}では描画されていた
顔の中の輪郭などが検出されていないことがわかる。


\subsection{フィルタ処理}

画像に対してローパスフィルタ(LPF)やハイパスフィルタ(HPF)による
フィルタリングをすることが出来る。
LPFではノイズの除去や画像のぼかしが、
HPFではエッジの強調ができる。

filter2D関数は、入力画像とカーネル(フィルタ)の畳み込みを計算する。
式\ref{fomula:kernel}は、画像の平滑化を表す。
\begin{eqnarray}
  K ~=~ \frac{1}{25}
  \left[
    \begin{array}{rrrrr}
      1&1&1&1&1\\
      1&1&1&1&1\\
      1&1&1&1&1\\
      1&1&1&1&1\\
      1&1&1&1&1
    \end{array}
  \right]
  \label{fomula:kernel}
\end{eqnarray}

式\ref{fomula:kernel}を使ったフィルタリングを行うと、
各画素に対してその画素を中心に$5\times5$の画素を選択し、
その合計を25で割ったものを出力画像でのその画素の値にすることが出来る。

OpenCVではエッジ検出が出来るCannyや、
エッジを劣化させずに平滑化できるbilatetalFilterなど
多くのフィルタが利用できる。


\subsection{自動検出によるマスク処理}

読み込んだ画像は通常、画素ごとに[G,B,R]のarrayでデータを持つ。
これに対して図\ref{fig:mask1}のようなマスク画像を用意し
bitwise\_andをとると、図\ref{fig:masked1}のような画像が生成できる。

\begin{figure}
  \begin{minipage}{0.5\hsize}
    \centering
    \includegraphics[width=0.9\hsize]{../pic/black.jpg}
    \caption{マスク画像}
    \label{fig:mask1}
  \end{minipage}
  \begin{minipage}{0.5\hsize}
    \centering
    \includegraphics[width=0.9\hsize]{../pic/tapu_icon.jpg}
    \caption{bitwise\_andを撮った画像}
    \label{fig:masked1}
  \end{minipage}
\end{figure}




\section{結果と考察}



\begin{thebibliography}{10}
  \bibitem{1} OpenCV-Python チュートリアル文書
  
  \url{http://labs.eecs.tottori-u.ac.jp/sd/Member/oyamada/OpenCV/html/index.html}
  
  \bibitem{2} OpenCV / findContours を使用して画像中のオブジェクトの輪郭を検出する方法
  
  \url{https://axa.biopapyrus.jp/ia/OpenCV/detect-contours.html}
\end{thebibliography}
\end{document}